% !TEX root = pennstandermath_guide.tex
\usetypescriptfile[pennstander]

\usemodule[math-characters-renewed]

\setupbodyfont[pennstander-regular]
%\setupbodyfont[pennstander-black]
%\setupbodyfont[pennstander-thin]

\definefontfamily[dejavu][rm][DejaVu Serif]
\switchtobodyfont[dejavu]

\setuppagenumbering[state=stop]

\setuphead
  [title]
  [style=\dejavu,align=middle]


\setuphead
  [section]
  [style=\dejavu,number=no,align=middle]

\title[title={Pennstander Math}]

\starttext

\startbuffer[weightdisplays]
\starttabulate[|c|l|]
\NR \NC Thin \NC {\switchtobodyfont[pennstander-thin]
 \dm{
\int_a^b \!\! f'(x) dx = f(b) -f(a)
}}

\switchtobodyfont[pennstander-extralight] 
\NR \NC Extralight \NC \dm{
c:= \sqrt{a^2 + b^2} \notin \mathbb Q
}

\NR \NC Light \NC {\switchtobodyfont[pennstander-light]\dm{
\oint_\gamma \hat{f}(z) dz =0
}}

\NR \NC Regular \NC {\switchtobodyfont[pennstander-regular]
\dm{
\iint_M d\omega = \int_{\partial M}\omega}\blank[10pt]
}

\NR \NC  Medium \NC
{\switchtobodyfont[pennstander-medium]
\dm{
\lim_{n\to \infty}\sum_{\mu=1}^n \frac{1}{\mu^2} = \frac{\pi^2}{6}
}\blank[10pt]}

\NR \NC 
 Semibold \NC { \switchtobodyfont[pennstander-semibold] \switchtobodyfont[pennstander-semibold]

\dm{
\frac{d}{d \tau}(\psi\cdot \phi) = \frac{d\psi}{d\tau} \cdot \phi + \psi \cdot \frac{d\phi}{d\tau}
}}


\NR \NC 
 Bold \NC {\switchtobodyfont[pennstander-bold]

\dm{
\nabla\cdot \fenced[group]{ \nabla \times G}= 0
}}

\NR \NC  ExtraBold \NC {\switchtobodyfont[pennstander-extrabold] \dm{\fenced[brace]{e^{i\theta} = \cos(\theta) + i \sin(\theta) : \theta\in\mathbb R}}\blank[10pt]
}


\NR \NC Black \NC
{\switchtobodyfont[pennstander-black] \dm{\Psi \otimes \Omega = \Omega \otimes \Psi}}
\stoptabulate
\stopbuffer

\startbuffer[weightdisplays]
\startalign[center]
\setupTABLE[frame=off, valign=middle, align=middle]
\setupTABLE[column][2][width=1cm]
\bTABLE
  \bTR[height=1cm]
    \bTD Thin \eTD
    \bTD \eTD
    \bTD Extralight \eTD
  \eTR
  \bTR[height=2cm]
    \bTD {\switchtobodyfont[pennstander-thin]
      \dm{\int_a^b f'(x) dx = f(b) -f(a)}} \eTD
    \bTD \eTD
    \bTD {\switchtobodyfont[pennstander-extralight]
      \dm{c:= \sqrt{a^2 + b^2} \notin \mathbb Q}} \eTD
  \eTR

  \bTR[height=1cm]
    \bTD Light \eTD
    \bTD \eTD
    \bTD Regular \eTD
  \eTR
  \bTR[height=2cm]
    \bTD {\switchtobodyfont[pennstander-light]
      \dm{\oint_\gamma \hat{f}(z) dz = 0}} \eTD
    \bTD \eTD
    \bTD {\switchtobodyfont[pennstander-regular]
      \dm{\iint_M d\omega = \int_{\partial M}\omega}} \eTD
  \eTR

  \bTR[height=1cm]
    \bTD Medium \eTD
    \bTD \eTD
    \bTD Semibold \eTD
  \eTR
  \bTR[height=2cm]
    \bTD {\switchtobodyfont[pennstander-medium]
      \dm{\lim_{n\to \infty}\sum_{\mu=1}^n \frac{1}{\mu^2} = \frac{\pi^2}{6}}} \eTD
    \bTD \eTD
    \bTD {\switchtobodyfont[pennstander-semibold]
      \dm{\frac{d}{d \tau}(\psi\cdot \phi) = \frac{d\psi}{d\tau} \cdot \phi + \psi \cdot \frac{d\phi}{d\tau}}} \eTD
  \eTR

  \bTR[height=1cm]
    \bTD Bold \eTD
    \bTD \eTD
    \bTD ExtraBold \eTD
  \eTR
  \bTR[height=2cm]
    \bTD {\switchtobodyfont[pennstander-bold]
      \dm{\nabla\cdot \fenced[group]{\nabla \times G} = 0}} \eTD
    \bTD \eTD
    \bTD {\switchtobodyfont[pennstander-extrabold]
      \dm{\fenced[brace]{e^{i\theta} = \cos(\theta) + i \sin(\theta) : \theta\in\mathbb R}}} \eTD
  \eTR

  \bTR[height=1cm]
    \bTD Black \eTD
    \bTD \eTD
    \bTD \eTD
  \eTR
  \bTR[height=2cm]
    \bTD {\switchtobodyfont[pennstander-black]
      \dm{\Psi \otimes \Omega = \Omega \otimes \Psi}} \eTD
    \bTD \eTD
    \bTD \eTD
  \eTR
\eTABLE
\stopalign
\stopbuffer
\getbuffer[weightdisplays]
\define\testcharacter{{\,\, }}

\define[1]\AllFenceSizes{
\im{
\fenced[#1][size=12]{
\fenced[#1][size=11]{
\fenced[#1][size=10]{
\fenced[#1][size=9]{
\fenced[#1][size=8]{
\fenced[#1][size=7]{
\fenced[#1][size=6]{
\fenced[#1][size=5]{
\fenced[#1][size=4]{
\fenced[#1][size=3]{
\fenced[#1][size=2]{
\fenced[#1][size=1]{
\fenced[#1][size=0]{\testcharacter}}}}}}}}}}}}}
}
}

\define\AllFences{
\AllFenceSizes{angle}
\AllFenceSizes{bracket}\crlf\crlf
\AllFenceSizes{brace}
\AllFenceSizes{bar}\crlf\crlf
\AllFenceSizes{parenthesis}
\AllFenceSizes{doublebar}\crlf\crlf
\AllFenceSizes{group}
\AllFenceSizes{ceiling}\crlf\crlf
\AllFenceSizes{floor}
\AllFenceSizes{triplebar}\crlf\crlf
\AllFenceSizes{openbracket}
}
\def\fontlist{pennstander-thin,pennstander-regular,pennstander-bold}
\def\fontlist{pennstander-regular}

\startbuffer[alphabets]
\doloopoverlist\fontlist{
{\switchtobodyfont[\recursestring]
   \doloopoverlist\alphabetlist
{%
\startalignment[center]
\showmathfontcharacters[alternative=glyphs,list=\recursestring,showfeatures=none]
\stopalignment
}
}}
\stopbuffer
\page
\setupmathematics[default=normal,lcgreek=normal]
\section[title={Upright Latin and Greek upper and lower case}]
\def\alphabetlist{lowercasenormal,uppercasenormal,digitsnormal,uppercasegreeknormal,lowercasegreeknormal}
\getbuffer[alphabets]
\section[title={Italic/Oblique Latin and Greek upper and lower case}]

\def\alphabetlist{lowercaseitalic,uppercaseitalic,uppercasegreekitalic,lowercasegreekitalic}
\getbuffer[alphabets]
\subsection[title={Double Struck upper case and numerals}]
\def\alphabetlist{uppercasedoublestruck,digitsdoublestruck}
\getbuffer[alphabets]
\subsection[title={Fraktur-like upper and lower case}]
\def\alphabetlist{uppercasefraktur,lowercasefraktur}
\getbuffer[alphabets]



\subsection[title={Script upper case}]
\def\alphabetlist{uppercasescript}
\getbuffer[alphabets]

\page
\subsection[title={Integrals}]

\startbuffer[integrals]
 \define\integrand{}
\define[1]\testintegrals{
\int[size=#1]{\integrand}  \iint[size=#1]{\integrand} \iiint[size=#1]{\integrand} 
\iiiint[size=#1]{\integrand} \oint[size=#1]{\integrand}  \oiint[size=#1]{\integrand}  \oiiint[size=#1]{\integrand}  \slashint[size=#1]{\integrand}   \barint[size=#1]{\integrand} \,\, 
 \aointc[size=#1]{\integrand} \,\, \ointc[size=#1]{\integrand} \,\, \intc[size=#1]{\integrand} \,\, \aodownintc[size=#1]{\integrand} \,\, \doublebarint[size=#1]{\integrand}\,\,\hookleftarrowint[size=#1]{\integrand} \,\,
 \timesint[size=#1]{\integrand} \,\, \capint[size=#1]{\integrand} \,\,  \cupint[size=#1]{\integrand} \,\,  \upperint[size=#1]{\integrand} \,\,  \lowerint[size=#1]{\integrand} \,\,
  \rectangularpoleintc[size=#1]{\integrand} \,\, \semicirclepoleintc[size=#1]{\integrand} \,\,
  \circlepoleoutsideintc[size=#1]{\integrand} \,\,\circlepoleinsideintc[size=#1]{\integrand} \,\, \squareintc[size=#1]{\integrand} \,\,
}
\startformula
\testintegrals{.15cm}\breakhere
\testintegrals{.4cm}\breakhere
\testintegrals{1cm}\breakhere
\stopformula
\stopbuffer

\doloopoverlist\fontlist
{
{\switchtobodyfont[\recursestring]
\getbuffer[integrals]}
}


\subsection[title={Radicals}]
\startbuffer[radicals]
\startformula
\sqrt[2]{1+x+y}+
\sqrt[1/3]{\frac{a+b}{c+d}}
\stopformula 
\stopbuffer


\doloopoverlist\fontlist
{
{\switchtobodyfont[\recursestring]
\getbuffer[radicals]}
}
\subsection[title={Accents and Stackers}]
\startbuffer[accents]
\startformula
\grave{x} \quad \tilde{x} \quad \dot{x}\quad \ddot{x}\quad \dddot{x}\quad\ddddot{x}\quad \check{x}\quad \acute{x} \quad\bar{x}\quad\breve{x}\quad\ring{x}
\stopformula
\startformula
\widehat{x+y}\quad
\widebar{x+y}\quad
\wideoverleftharpoon{x+y}\quad
\wideoverleftrightarrow{x+y}\quad
\wideunderleftrightarrow{x+y}\quad
\wideunderleftharpoon{x+y}\quad
\wideunderrightarrow{x+y}\quad
\widecheck{x+y}\quad\widetilde{x+y}\quad
\breakhere
%\wideoverrightharpoon{x+y}\quad
%\wideoverrightarrow{x+y}\quad
%\wideunderbar{x+y}\quad
%\wideunderrightharpoon{x+y}\quad
%\wideunderleftarrow{x+y}\quad
\overbrace{x+y}\quad
\overparent{x+y}\quad
\overbracket{x+y}\quad
\overbar{x+y}\quad
\breakhere
\underbrace{x+y}\quad
\underparent{x+y}\quad
\underbracket{x+y}\quad
\underbar{x+y}\quad
\stopformula
\stopbuffer

\doloopoverlist\fontlist
{
{\switchtobodyfont[\recursestring]
\getbuffer[accents]}
}


\subsection[title={Arrows}]
\startbuffer[arrows]
\startformula
A \mrel{1+2}{a+b+c}  B \mequal{1+2}{a+b+c} 
C \mleftarrow{1+2}{a+b+c} D \mrightarrow{1+2}{a+b+c} 
E \mleftrightarrow{1+2}{a+b+c} F \mLeftarrow{1+2}{a+b+c}
G \mRightarrow{1+2}{a+b+c} H \mLeftrightarrow{1+2}{a+b+c}
I \mtwoheadleftarrow{1+2}{a+b+c} J \mtwoheadrightarrow{1+2}{a+b+c}
K \mmapsto{1+2}{a+b+c} L \mhookleftarrow{1+2}{a+b+c}
M \mhookrightarrow{1+2}{a+b+c} N \mleftharpoondown{1+2}{a+b+c}
O \mleftharpoonup{1+2}{a+b+c} P \mrightharpoondown{1+2}{a+b+c}
Q \mrightharpoonup{1+2}{a+b+c} R \mrightoverleftarrow{1+2}{a+b+c}
S \mleftoverrightarrow{1+2}{a+b+c} T \mleftrightharpoons{1+2}{a+b+c}
U \mrightleftharpoons{1+2}{a+b+c} V \mtriplerel{1+2}{a+b+c} W
\stopformula
\stopbuffer


{\setupbodyfont[pennstander-regular]
\getbuffer[arrows]}

\subsection[title={Symbols}]

\def\symbollist{
    arrows,
%    basiclatin,
    blockelements,
%    combiningdiacriticalmarks,
%    combiningdiacriticalmarksforsymbols,
    controlpictures,
    currencysymbols,
    dingbats,
%    generalpunctuation,
    geometricshapes,
    ipaextensions,
    latinextendeda,
    latinsupplement,
    letterlikesymbols,
    mathematicaloperators,
    miscellaneousmathematicalsymbolsa,
    miscellaneousmathematicalsymbolsb,
    miscellaneoussymbols,
    miscellaneoussymbolsandarrows,
    miscellaneoustechnical,
    phoneticextensions,
    superscriptsandsubscripts,
    supplementalarrowsa,
    supplementalarrowsb,
    supplementalarrowsc,
    supplementalmathematicaloperators,
    supplementalsymbolsandpictographs,
    emoticons,
    miscellaneoussymbolsandpictographs}


\doloopoverlist\fontlist{
\edef\fontchoice{\recursestring}%
\doloopoverlist\symbollist{%
{\switchtobodyfont[\fontchoice]{\showmathfontcharacters[alternative=glyphswithoutsizes,list=\recursestring,showfeatures=none]}}}\blank[10pt]
}








\subsection[title=Brackets]
\startbuffer[brackets]
{\AllFences \blank[20pt]}
\stopbuffer

{\switchtobodyfont[pennstander-regular]
\getbuffer[brackets]}





\def\alternativeslist{
%    uppercasenormal,
%    lowercasenormal,
%    digitsnormal,
    digitsdoublestruck,
    uppercasefraktur,
    lowercasefraktur,
    arrows,
    basiclatin,
    blockelements,
%    combiningdiacriticalmarks,
%    combiningdiacriticalmarksforsymbols,
    controlpictures,
    currencysymbols,
    dingbats,
%    generalpunctuation,
    geometricshapes,
    ipaextensions,
    latinextendeda,
    latinsupplement,
    letterlikesymbols,
    mathematicaloperators,
    miscellaneousmathematicalsymbolsa,
    miscellaneousmathematicalsymbolsb,
    miscellaneoussymbols,
    miscellaneoussymbolsandarrows,
    miscellaneoustechnical,
    phoneticextensions,
    superscriptsandsubscripts,
    supplementalarrowsa,
    supplementalarrowsb,
    supplementalarrowsc,
    supplementalmathematicaloperators,
    supplementalsymbolsandpictographs}


\setupmathematics[default=normal]
{\subsection[title=Alternatives]
\switchtobodyfont[pennstander-regular]
\doloopoverlist{ss01,ss02,ss03,ss04,ss05,ss07,ss20,cv01,cv10,cv20,cv21,cv30,cv31,cv32,zero}{
{\switchtobodyfont[dejavu] \recursestring : }
\setupbodyfont[pennstander-regular]\showmathfontcharacters[alternative=glyphvariants,list=\alternativeslist,showfeatures=none,togglelist={\recursestring}]\crlf
}
}





%\doloopoverlist{pennstander-thin-ss,pennstander-regular-ss,pennstander-bold-ss}
%{
%\setupbodyfont[\recursestring]
%\getbuffer[alternatives]
%}




\subsection[title=Private Space Characters]
\startbuffer[additionals]
\starttabulate[|l|l|]
\NC U+E1001 \NC {\switchtobodyfont[pennstander-regular]$\char"E1001$}\NR
\NC U+E1002 \NC {\switchtobodyfont[pennstander-regular]$\char"E1002$}\NR
\NC U+E1003 \NC {\switchtobodyfont[pennstander-regular]$\char"E1003$}\NR
\NC U+E1004 \NC {\switchtobodyfont[pennstander-regular]$\char"E1004$}\NR
\NC U+E1005 \NC {\switchtobodyfont[pennstander-regular]$\char"E1005$}\NR
\NC U+E1006 \NC {\switchtobodyfont[pennstander-regular]$\char"E1006$}\NR
\NC U+E1007 \NC {\switchtobodyfont[pennstander-regular]$\char"E1007$}\NR
\stoptabulate
\stopbuffer
\switchtobodyfont[dejavu] 
\getbuffer[additionals]

\setupbodyfont[pennstander-bold]
\startformula
{\char"E9999}
\stopformula

% \startformula
% \rm agIAMNQVWZX \breakhere
% \it agIAMNQVWZX \breakhere
% \bf agIAMNQVWZX \breakhere
% \bi agIAMNQVWZX
% \stopformula
    

\stopcomponent